% Lösungen Einheit 1 von 2 – Jahreszins, Monats- und Tageszins (zusammenbau v0.1)
\einheitenkopf{Einheit 1 von 2 · Jahreszins, Monats- und Tageszins}

\erg{11}{a) die Zinsen (Prozentwert) \quad b) 6\,€ ($300 \cdot 0{,}02$) \quad c) 24\,€ (1\,\% = 6\,€, 4\,\% = 24\,€) \quad d) 714\,€ (Zinsen: $700 \cdot 0{,}02 = 14$\,€) \quad e) 3\,\% ($27 : 900 = 0{,}03$) \quad f) 1,5\,\% ($60 : 4\,000 = 0{,}015$) \quad g) 700\,€ (1\,\% = 7\,€) \quad h) $300 \cdot 0{,}03 = 9{,}00$\,€ – stimmt (oder $9 : 300 = 0{,}03 = 3\,\%$). \quad i) 36\,€ (Jahreszinsen 12\,€, $12 \cdot 3 = 36$) \quad j) 1\,908\,€ (Zinsen: $1\,800 \cdot 0{,}06 = 108$\,€) \quad k) Bank B: 65\,€, Bank A: 60\,€. \quad l) 52\,€ ($2\,600 \cdot 0{,}02$)}
\erg{12}{a) Monate – durch zwölf teilen \quad b) 20\,€ (Jahreszinsen 48\,€, $48 : 12 \cdot 5 = 20$) \quad c) 16\,€ (Jahreszinsen 144\,€, $144 : 360 \cdot 40 = 16$) \quad d) 45\,€ (75 Tage; Jahreszinsen 216\,€, $216 : 360 \cdot 75 = 45$) \quad e) 2\,\% (Jahreszinsen $20 : 4 \cdot 12 = 60$\,€, $60 : 3\,000 = 0{,}02$) \quad f) 52,50\,€ (Jahreszinsen 105\,€, $105 : 12 \cdot 6 = 52{,}50$). Nora bekommt 52,50\,€ Zinsen.}
\erg{13}{a) ≈ 80\,€ (1\,\% ≈ 40\,€, mal 2)}
\erg{14}{a) Gefragt sind die Zinsen, nicht das Guthaben: richtig 38\,€. \quad b) Schon 1\,\% von 20\,000\,€ sind 200\,€, 4\,\% von 600\,€ nur 24\,€ – der kleine Anteil vom viel größeren Kapital ist mehr. \quad c) Nein: Rückzahlung 4\,240\,€, gespart 4\,200\,€, es fehlen 40\,€.}
